\documentclass{ltjsarticle}
\usepackage{graphicx}
\usepackage{color}
\usepackage{url}
\usepackage{listings}
\usepackage{booktabs}

\begin{document}
\pagestyle{empty}

\begin{center}
\vspace*{1cm}
\large
{\LARGE 卒業研究報告書}\\
\vspace*{0.8cm}
題目\\
\vspace*{1cm}
{\Huge {マルチモーダル学習による薬剤鑑別}}\\
\vspace{3mm}

\vspace*{3cm}
指導教員\\
\vspace*{0.3cm}
{\LARGE 森山 真光  准教授}\\
\vspace*{3cm}
報告者\\
\vspace*{0.3cm}

{23--1--211--0009}\\
\vspace*{0.3cm}
{\Huge 呉竹 櫂人}\\
\vspace*{0.5cm}
近畿大学情報学部情報学科\\
\vspace*{2cm}
令和7年5月4日提出\\
\end{center}

\newpage
\normalsize

\clearpage
\section*{概要}

薬剤の取り違えや誤服用は医療安全上の重要な課題であり，医療現場や家庭では，薬剤名が記載された外箱や薬袋が失われた場合，または薬剤本体やPTPシートのみが残された場合に，錠剤やカプセルの外観，刻印，包装上の印字など限られた視覚情報を手がかりに薬剤を鑑別する必要がある．本研究では，薬剤画像とテキストプロンプトを入力とし，薬剤名を出力するVision-Language Model型の薬剤鑑別手法を検討する．具体的には，オープンソースのマルチモーダルLLMであるQwen3.5を用い，薬剤画像に対して「この錠剤は何という薬ですか」という質問を与え，正解薬剤名を出力するようにLoRAによるファインチューニングを行う．比較対象として，ファインチューニング前のQwen3.5およびChatGPT・Geminiなどの外部API型マルチモーダルAIによるゼロショット識別を用い，薬剤画像データに特化したファインチューニングが薬剤鑑別精度に与える効果を評価する．また，今後の発展的課題として，薬剤データベースを検索して候補薬剤情報を提示するRAG型手法との比較も検討する．

\newpage

\tableofcontents
\newpage
\setcounter{page}{1}
\pagestyle{plain}

%----------------------------------------------------------------------
\section{序論}\label{sec1}
%----------------------------------------------------------------------

\subsection{本研究の背景}

薬剤の取り違えや誤服用は，患者の健康被害につながる医療安全上の重要な問題である．
医療現場や家庭では，薬剤名が記載された外箱や薬袋が失われた場合，または薬剤本体やPTPシートのみが残された場合に，錠剤やカプセルの外観，刻印，包装上の印字など限られた視覚情報を手がかりに薬剤を鑑別する必要がある．
特に，高齢者が複数の薬剤を長期的に服用する場合や，一包化調剤によって複数の薬剤がまとめられる場合には，薬剤の識別が困難になり，誤服用のリスクが高まる．

従来，薬剤鑑別には薬剤師などの専門知識が必要であり，薬剤データベースを手作業で検索する負担も大きい．
また，白色円形の錠剤のように外観が類似する薬剤も多く，画像のみから薬剤を正確に識別することは容易ではない．
このような背景から，薬剤画像を用いて自動的に薬剤を識別する研究が行われてきた．

\subsection{従来研究と課題}

Heoら\cite{heo2023pill}は，薬剤識別を認識ステップと検索ステップに分けた深層学習ベースの薬剤識別システムを提案している．
YOLOv5により錠剤表面の刻印文字を検出し，ResNet-32により形状，色，剤形を認識する．
さらに，RNNベースの文字レベル言語モデルを用いて検出された刻印文字の誤りを補正し，薬剤データベース内の情報との類似度に基づいて候補薬剤を順位付けする．
この手法は，薬剤画像に含まれる視覚情報と刻印文字情報を組み合わせて識別する点で有効であり，検索型の構成を採用することで学習に含まれない薬剤への対応も考慮している．
一方で，複数の専用モジュールを組み合わせる必要があり，システム構成が複雑である．
また，刻印文字を重要な識別手がかりとしているため，刻印が不明瞭な画像や刻印のない薬剤では識別性能が低下する可能性がある．

Rádliら\cite{radli2024metric}は，薬剤画像の視覚特徴と薬剤説明テキストから得られるテキスト特徴を組み合わせたマルチモーダルな薬剤認識手法を提案している．
この研究は，画像情報だけでなくテキスト情報を利用して薬剤認識を行う点で本研究と近い．
一方で，同手法は画像特徴抽出器，テキスト埋め込み，メトリクス学習を個別に設計する構成であり，画像とテキストプロンプトを入力として薬剤名を直接生成するVision-Language Model型の手法ではない．

Menahaら\cite{gemini2024pill}は，Google Geminiを用いた生成AIベースの薬剤識別支援システムを提案している．
この研究は，マルチモーダルAIを薬剤識別支援に応用している点で本研究と関連する．
一方で，同研究は商用クローズドソースモデルを用いたプロンプトベースの手法であり，薬剤画像データを用いてモデル自体をファインチューニングする構成ではない．
また，外部APIを利用する場合，薬剤画像を外部サービスに送信する必要があり，プライバシー上の懸念が残る．

以上の課題から，（1）錠剤本体の外観，刻印，包装上の印字など複数の視覚情報を統合した識別，（2）外部APIに過度に依存しないローカル動作の可能性，（3）薬剤画像に特化したファインチューニングによる識別精度の向上，の3点を検討する必要がある．

\subsection{本研究の目的}

本研究では，オープンソースのVision-Language ModelであるQwen3.5を用い，薬剤画像とテキストプロンプトから薬剤名を出力する薬剤鑑別手法を検討する．
具体的には，薬剤画像に対して「この錠剤は何という薬ですか」という質問を与え，正解薬剤名を出力するようにLoRAによるファインチューニングを行う．

本研究の目的は，ファインチューニング前のQwen3.5およびChatGPT・Geminiなどの外部API型マルチモーダルAIと比較することで，薬剤画像データに特化したファインチューニングが薬剤鑑別精度に与える効果を検証することである．
また，オープンソースモデルを用いることで，薬剤画像を外部APIに送信せず，適切な計算環境を用意すればローカル環境で推論できる薬剤鑑別システムの可能性を示すことを目指す．

なお，本研究におけるマルチモーダル学習とは，薬剤画像を視覚情報として入力し，自然言語による質問文をテキスト情報として与え，薬剤名を応答として生成するVision-Language Model型の学習を指す．
本報告時点では，テキスト入力は質問文を中心とするが，今後は識別コード，薬剤候補名，形状情報，効能情報などを追加し，より情報量の多い画像・テキスト統合型の鑑別へ拡張することを検討する．

\subsection{本報告書の構成}

第2章では関連研究を整理し，本研究の位置づけを明確にする．
第3章では提案手法の詳細および実験計画を述べる．
第4章では実験結果と考察を示す．
第5章では結論と今後の課題を述べる．

\newpage

%----------------------------------------------------------------------
\section{関連研究}\label{sec2}
%----------------------------------------------------------------------

本研究では以下の3本の先行研究を参考文献とし，各手法の概要と限界を整理する．

\subsection{画像特徴・刻印文字を用いた薬剤識別手法}

Heoら\cite{heo2023pill}は，薬剤識別を認識ステップと検索ステップに分けた深層学習ベースの薬剤識別システムを提案した．
認識ステップでは，YOLOv5により錠剤表面の刻印文字を検出し，ResNet-32により形状，色，剤形を認識する．
さらに，RNNベースの文字レベル言語モデルを用いて，検出された刻印文字の誤りを補正する．
検索ステップでは，認識された形状，色，剤形，刻印文字と薬剤データベース内の情報との類似度を計算し，候補薬剤を順位付けする．
韓国の薬剤データセット（MFDS）においてTop-1精度85.6\%，米国データセット（NLM）において74.5\%を達成している．

この手法は，薬剤画像に含まれる視覚情報と刻印文字情報を組み合わせて識別する点で有効であり，画像分類ではなく検索型の構成を採用することで，学習に含まれない薬剤への対応も考慮している．

\textbf{課題}
\begin{itemize}
  \item YOLOv5，ResNet，RNNベースの言語モデルなど複数のモジュールを組み合わせる必要があり，システム構成が複雑である．
  \item 刻印文字を重要な識別手がかりとしているため，刻印が不明瞭な画像や刻印のない薬剤では，識別性能が低下する可能性がある．
  \item 検索型の構成により未知薬剤への対応を考慮している一方で，刻印検出や特徴認識の各モジュールが対象データに依存する可能性がある．
\end{itemize}

\subsection{画像特徴とテキスト特徴を用いたマルチモーダル薬剤識別手法}

Rádliら\cite{radli2024metric}は，薬剤画像の視覚特徴と薬剤説明文から得られるテキスト特徴を組み合わせたマルチモーダルな薬剤認識手法を提案している．
同研究では，EfficientNetV2をバックボーンとし，RGB画像，局所特徴，輪郭，テクスチャなど複数の視覚特徴を抽出する．
また，薬剤説明文から得られる単語埋め込みを用いて，メトリクス学習における動的マージントリプレット損失（DMTL）のマージンを決定する．
112クラス・4,480枚で構成するOGYEIv2データセットにてTop-1精度93.56\%を達成している．

この研究は，画像情報だけでなくテキスト情報を利用して薬剤認識を行う点で本研究と近い．

\textbf{課題}
\begin{itemize}
  \item 画像特徴抽出器，テキスト埋め込み，メトリクス学習を個別に設計する必要がある．
  \item 画像とテキストプロンプトから薬剤名を直接出力するVLM型の手法ではない．
  \item 評価データが特定地域・特定データセットに限定されており，他の薬剤画像への適用可能性は追加検証が必要である．
\end{itemize}

\subsection{生成AI・LLMを用いた薬剤識別支援手法}

Menahaら\cite{gemini2024pill}は，Google Geminiを用いた生成AIベースの薬剤識別支援システムを提案している．
同研究は，薬剤画像を入力として，薬剤に関する情報をテキストおよび音声で提示するシステムであり，マルチモーダルAIを薬剤識別支援に応用している点で本研究と関連する．

一方で，同研究は商用クローズドソースモデルを用いたプロンプトベースの手法である．
少なくとも同研究では，薬剤画像データを用いてモデル自体をファインチューニングする構成は示されていない．
また，外部APIを利用する場合，薬剤画像を外部サービスに送信する必要があり，プライバシー上の懸念が残る．
本研究では，オープンソースモデルであるQwen3.5を用い，ローカル環境でLoRAファインチューニングを行う点に差異がある．

\textbf{課題}
\begin{itemize}
  \item 商用クローズドソースモデルであり，当該研究ではモデル自体のファインチューニングは行われていない．
  \item 外部APIを利用する場合，薬剤画像を外部サービスに送信する必要があり，プライバシー上の懸念が残る．
  \item 利用コストが高く，継続的な運用において制約が生じる可能性がある．
  \item プロンプトベースの手法であり，当該研究では対象薬剤への追加学習は示されていない．
\end{itemize}

\subsection{先行研究の整理と本研究の位置づけ}

先行研究では，薬剤画像のみを用いるCNN分類に加え，刻印文字，薬剤説明テキスト，処方情報などを組み合わせるマルチモーダルな薬剤識別手法が提案されている．
Heoらの手法は，画像特徴と刻印文字を個別に抽出し，言語モデルで刻印文字を補正することで高精度な薬剤検索を実現している．
一方で，複数の専用モジュールを組み合わせる必要があり，システム構成が複雑である．
Rádliらの手法は，薬剤画像と薬剤説明テキストを組み合わせる点で本研究と関連するが，画像特徴抽出器とテキスト埋め込みを用いたメトリクス学習であり，薬剤名を直接生成するVLM型の手法ではない．
Menahaらの手法は，生成AIを薬剤識別支援に用いる点で関連するが，商用モデルを用いたプロンプトベースの手法であり，薬剤画像に特化したファインチューニングを行うものではない．

これらを踏まえ，本研究ではオープンソースのVision-Language ModelであるQwen3.5を用い，薬剤画像とテキストプロンプトを入力として薬剤名を直接出力する手法を検討する．
特に，LoRAによるファインチューニングを行うことで，薬剤画像データに特化したモデルを構築し，ファインチューニング前のモデルおよび外部API型マルチモーダルAIとの比較を通じて，その有効性を評価する．

表\ref{tab:related}に先行研究と本研究の比較を示す．

\begin{table}[h]
\centering
\caption{先行研究と本研究の位置づけ}
\label{tab:related}
\small
\setlength{\tabcolsep}{3pt}
\begin{tabular}{p{0.16\linewidth}p{0.10\linewidth}p{0.14\linewidth}p{0.24\linewidth}p{0.28\linewidth}}
\hline
手法 & 画像利用 & テキスト利用 & プライバシー & 本研究との差異 \\
\hline
Heoら \cite{heo2023pill}     & あり & 刻印文字 & 低リスク & 複数モジュール構成，VLM型ではない \\
Rádliら \cite{radli2024metric} & あり & 薬剤説明文 & 低リスク & メトリクス学習，薬剤名の直接生成ではない \\
Menahaら \cite{gemini2024pill} & あり & プロンプト & 外部API利用時に懸念 & 商用モデル，当該研究ではFTなし \\
\textbf{本研究}               & あり & 質問文 & 外部APIに依存しない構成を目指す & Qwen3.5をLoRAで薬剤画像に適応 \\
\hline
\end{tabular}
\end{table}

\newpage

%----------------------------------------------------------------------
\section{研究内容}\label{sec3}
%----------------------------------------------------------------------

\subsection{使用モデル：Qwen3.5}

本研究では，画像とテキストを入力できるオープンソースのVision-Language ModelとしてQwen3.5を使用する．
Qwen3.5を選定する理由は，外部APIに依存せずローカル環境で動作させられる可能性があり，LoRAなどのパラメータ効率の高い手法によって追加学習が可能であるためである．

Vision-Language Modelは画像とテキストプロンプトを同時に入力できるため，薬剤画像に対する質問応答形式で薬剤名を出力するタスクを構成できる．
本研究では，この特性を利用し，薬剤画像から薬剤名を出力するようにQwen3.5をLoRAによりファインチューニングする．

\subsection{提案手法：Qwen3.5 の LoRA ファインチューニング}

本研究では，薬剤画像とテキストプロンプトを入力し，正解薬剤名を出力するようにQwen3.5をLoRAによりファインチューニングする．
学習データは，薬剤ごとに分類された画像フォルダから作成する．
各サブフォルダ名を薬剤名ラベルとし，フォルダ内の画像を入力画像として用いる．

入力プロンプトは「この錠剤は何という薬ですか。薬剤名だけで答えてください。」とし，モデルが正解薬剤名のみを出力するように学習させる．
これにより，薬剤画像に対する質問応答形式の薬剤鑑別モデルを構築する．

\textbf{強み}
\begin{itemize}
  \item 薬剤画像データに特化した追加学習により，ゼロショット識別よりも高い識別精度が期待できる．
  \item 外部APIに依存せず，適切な計算環境を用意すればローカル環境で推論できる．
  \item 画像と自然言語プロンプトを組み合わせた質問応答形式で薬剤名を出力できる．
\end{itemize}

\textbf{課題}
\begin{itemize}
  \item ある程度の量と多様性を持つ学習データが必要である．
  \item 新たな薬剤クラスを追加する場合，再チューニングまたはデータベース検索との併用が必要となる．
  \item 包装込み画像を用いる場合，モデルが錠剤本体ではなく包装上の文字やデザインに依存する可能性がある．
\end{itemize}

\subsection{比較手法}

提案手法の有効性を検証するため，以下の比較手法を設定する．

\begin{itemize}
  \item \textbf{ファインチューニング前のQwen3.5}：LoRA適用前のモデルによるゼロショット識別．
  \item \textbf{ChatGPT}：外部API型マルチモーダルAIによるゼロショット識別．
  \item \textbf{Gemini}：外部API型マルチモーダルAIによるゼロショット識別．
\end{itemize}

\subsection{使用データ}

本研究では，研究用に撮影した薬剤画像を使用する．
画像には，錠剤やカプセル本体に加え，PTPシートや包装上の印字が含まれる場合がある．
これは，実際の薬剤鑑別において，錠剤本体の色や形状だけでなく，包装上の文字情報や印字も重要な手がかりとなるためである．

ただし，包装情報を含む画像を用いる場合，モデルが錠剤本体ではなく包装上の文字やデザインに依存して薬剤名を予測する可能性がある．
そのため，評価時には，包装込み画像と錠剤部分を中心とした画像で性能差を比較し，モデルがどの情報に依存しているかを分析する必要がある．
本研究では，包装情報を含む画像をそのまま用いる場合と，錠剤またはカプセル部分を中心に切り出した画像を用いる場合を区別して評価することを検討する．
これにより，モデルが包装上の印字に依存しているのか，薬剤本体の色や形状などの外観特徴を利用しているのかを分析する．

\subsection{評価方法}

評価においては，学習画像と評価画像が重複しないようにデータを分割する．
具体的には，各薬剤クラスごとに画像を学習用と評価用に分け，同一画像および極端に類似した連続撮影画像が学習データと評価データの両方に含まれないようにする．
これにより，単なる画像の記憶ではなく，薬剤画像に対する識別性能を評価する．
評価では，入力画像に対してモデルが出力した薬剤名と正解薬剤名を比較し，識別精度を算出する．
主な評価指標として，Top-1 Accuracyを用いる．
Top-1 Accuracyは，モデルが最初に出力した薬剤名が正解薬剤名と一致した割合である．
また，モデルに上位3候補を出力させる実験を行う場合には，Top-3 Accuracyも補助指標として用いる．
Top-3 Accuracyを算出する場合は，「可能性の高い薬剤名を3件出力してください」という固定プロンプトを用い，上位3候補の中に正解薬剤名が含まれるかを判定する．
また，包装込み画像と錠剤部分を中心とした画像を分けて評価することで，包装印字への依存度を分析する．

さらに，誤識別例については，外観が類似している薬剤，包装印字が不明瞭な薬剤，撮影条件が悪い画像などに分類し，誤識別の原因を分析する．

\subsection{実験計画}

本研究では，以下の手順で実験を進める．

\begin{enumerate}
  \item 薬剤画像データセットを作成し，薬剤名ラベルを整理する．
  \item 各薬剤クラスごとに，画像を学習用，検証用，評価用に分割する．この際，同一画像および極端に類似した連続撮影画像が学習データと評価データの両方に含まれないようにする．
  \item ファインチューニング前のQwen3.5を用いて，薬剤画像に対するゼロショット識別を行う．
  \item 薬剤画像と正解薬剤名のペアを用いて，Qwen3.5をLoRAによりファインチューニングする．
  \item ファインチューニング後のモデルを用いて同一条件で識別実験を行う．
  \item ChatGPT・Geminiなどの外部API型マルチモーダルAIと比較する．
  \item 包装込み画像と錠剤部分を中心とした画像で性能差を比較し，包装印字への依存度を分析する．
  \item 誤識別例を分析し，薬剤外観，包装印字，撮影条件などが識別精度に与える影響を考察する．
\end{enumerate}

表\ref{tab:schedule}に実験スケジュールを示す．

\begin{table}[h]
\centering
\caption{実験スケジュール}
\label{tab:schedule}
\begin{tabular}{lll}
\hline
フェーズ & 時期 & 内容 \\
\hline
Phase 1 & 4〜5月  & 先行研究調査，研究課題と提案手法の整理 \\
Phase 2 & 5〜6月  & 薬剤画像データセットの整理，評価方法の確定 \\
Phase 3 & 6〜8月  & Qwen3.5のゼロショット評価とLoRAファインチューニング \\
Phase 4 & 8〜10月 & ChatGPT・Geminiとの比較，誤識別分析 \\
Phase 5 & 10〜1月 & 結果考察，卒業論文執筆，発表準備 \\
\hline
\end{tabular}
\end{table}

\newpage

%----------------------------------------------------------------------
\section{結果・考察}\label{sec4}
%----------------------------------------------------------------------

（実験実施後に記述予定）

\newpage

%----------------------------------------------------------------------
\section{結論・今後の課題}\label{sec5}
%----------------------------------------------------------------------

（実験実施後に記述予定）

\newpage

\nocite{*}
\bibliographystyle{plain}
\bibliography{reference}

\end{document}
